\documentclass[11pt,a4paper]{article}
\usepackage{fontspec}
\usepackage{unicode-math}
\setmainfont{STIX Two Text}
\setmathfont{STIX Two Math}
\usepackage[margin=2.2cm]{geometry}
\usepackage{booktabs,longtable,array}
\usepackage{xcolor}
\usepackage[hidelinks]{hyperref}
\usepackage{amsthm}
\newtheorem*{theorem}{Theorem}
\newcommand{\leanus}{\textunderscore\allowbreak}
\newcommand{\lean}[1]{{\small\ttfamily\let\_\leanus #1}}
\sloppy
\setlength{\parskip}{0.45em}
\setlength{\parindent}{0pt}

\title{NRS for every conjugate pair\\[0.3em]
\large The Robertson--Schrödinger table on $T_d{:}P_d$ at maximal tension}
\author{Eduardo Nava Hernández}
\date{26 September 2026 \\ \small Lean 4 package \texttt{nava-robertson-schrodinger}, module \texttt{D39\_ConjugatePairs}}

\begin{document}
\maketitle

\section{What the table says}

Robertson (1929) and Schrödinger (1930) proved, for any two observables $A$, $B$ and any
normalized state $\psi$,
\[
  \sigma_A^2\,\sigma_B^2 \;\ge\; \Big(\tfrac12\langle\{A,B\}\rangle-\langle A\rangle\langle B\rangle\Big)^2
  + \Big(\tfrac12\,|\langle[A,B]\rangle|\Big)^2 \;=\; |\langle \tilde A\psi,\tilde B\psi\rangle|^2 ,
\]
where $\tilde A\psi=(A-\langle A\rangle)\psi$ and $\tilde B\psi=(B-\langle B\rangle)\psi$ are the
fluctuation vectors. The theorem carries no list of pairs: it holds for every pair
(\lean{D2\_Robertson}). NRS evaluates it on the elemental pair $T_d{:}P_d$ at the state of
maximal tension $\psi^*$, and measures the ratio
\[
  R \;=\; \frac{\sigma_A\,\sigma_B}{|\langle \tilde A\psi^*,\tilde B\psi^*\rangle|}\;\ge 1 .
\]
$R=1$ means the inequality is saturated; $R>1$ is what the inequality cannot close.

\textbf{The result of this note.} $R$ is the same number $C_{\mathrm{Nava}}(d)$ for
\emph{every} conjugate pair realized on $T_d{:}P_d$. It equals $1$ exactly at $d=2,3$; from
$d=4$ on it is strictly larger than $1$, grows strictly with $d$, and stays strictly below
$C_\infty=\sqrt{\pi^2/3-2}$, which no $d$ reaches. On a cube of three pairs the leftover is
the volumetric quantum $\mathcal V$: zero at $3\times3\times3$, strictly positive from
$4\times4\times4$, growing with every axis, never vanishing. All of this is proved in Lean 4
with only the standard axioms.

\section{The setting}

On $H_d=\mathbb C^d$:
\begin{itemize}
  \item $P_d$ is diagonal with the $d$ equispaced values of $[-1,1]$: it orders the values of
        the pair.
  \item $T_d=A_d/\rho_d$ is the adjacency of the path on $d$ sites, $\rho_d=2\cos(\pi/(d+1))$:
        the step between consecutive values. It is the only graph compatible with ordered
        locality and completeness (\lean{D3\_PathGraph}).
  \item $\psi^*_j\propto(-i)^j\sin\big((j+1)\pi/(d+1)\big)$ is the maximal-tension state
        (\lean{D5\_MaximalTension}, \lean{D21}); at $\psi^*$,
        $\tfrac12|\langle[T_d,P_d]\rangle| = 1/(d-1)$.
\end{itemize}

\textbf{A pair realized on $T_d{:}P_d$} is
\[
  A = a\,T_d + b, \qquad B = c\,P_d + e, \qquad a,c\in\mathbb R\setminus\{0\},\; b,e\in\mathbb R .
\]
The numbers $a,c$ are the units of the pair and $b,e$ its origins: metres and kg\,m/s for
position--momentum, a count and radians for number--phase, coulombs and webers for
charge--flux, and so on.

\section{The theorem}

\begin{theorem}[\lean{D39\_ConjugatePairs}]
Let $d\ge2$ and $a,c\neq0$. For the pair $A=aT_d+b$, $B=cP_d+e$ at $\psi^*$:
\begin{enumerate}
  \item $R = C_{\mathrm{Nava}}(d)$ \hfill(\lean{razon\_par});
  \item the fluctuation vectors are parallel ($R=1$) if and only if $d=2$ or $d=3$
        \hfill(\lean{satura\_par\_iff});
  \item for $d\ge4$ they meet at an angle $\theta_{\mathrm{NRS}}(d)=\arccos(1/C_{\mathrm{Nava}}(d))>0$
        \hfill(\lean{angulo\_par\_pos});
  \item the angle grows strictly with $d$, even if the pair at $d$ and the pair at $d'$ carry
        different units \hfill(\lean{angulo\_par\_lt\_of\_lt});
  \item it stays strictly below $\arccos(1/C_\infty)\approx28.30^\circ$, which no $d$ attains
        \hfill(\lean{angulo\_par\_lt\_limite});
  \item the excess $R-1$ is the dimensional quantum $\delta(d)=C_{\mathrm{Nava}}(d)-1$
        \hfill(\lean{cuanto\_par}).
\end{enumerate}
On the cube $d_x\times d_y\times d_z$ with one pair per axis, $A_i=a_iT_{d_i}+b_i$,
$B_i=c_iP_{d_i}+e_i$ (all $a_i,c_i\neq0$, each axis in its own units) and the state
$\Psi^*=\psi^*\otimes\psi^*\otimes\psi^*$:
\begin{enumerate}\setcounter{enumi}{6}
  \item axis $i$ has ratio $C_{\mathrm{Nava}}(d_i)$ \hfill(\lean{razon\_par\_eje\_x/y/z});
  \item $(R_x-1)(R_y-1)(R_z-1)=\mathcal V(d_x,d_y,d_z)=\delta(d_x)\,\delta(d_y)\,\delta(d_z)$
        \hfill(\lean{cuantoVolumetrico\_pares});
  \item $\mathcal V=0$ if and only if some axis has $2$ or $3$ sites; from $4\times4\times4$ on,
        $0<\mathcal V<\delta_\infty^3$ \hfill(\lean{cuanto\_volumetrico\_pares}).
\end{enumerate}
\end{theorem}

\textbf{Why units do not matter.} For a unit state,
$\langle aL+b\rangle = a\langle L\rangle+b$, so the fluctuation vector of $aL+b$ is exactly
$a$ times that of $L$ (\lean{centradoG\_afin}). In
$R=\|\tilde A\psi\|\,\|\tilde B\psi\|/|\langle\tilde A\psi,\tilde B\psi\rangle|$ the factor
$|a|\,|c|$ appears once in the numerator and once in the denominator and cancels
(\lean{razonG\_afin}); the same holds for the angle (\lean{anguloG\_afin}). These two lemmas
hold at \emph{every} unit state, not only at $\psi^*$, so any statement about $T_d{:}P_d$
that depends only on $R$ or on the angle, such as strictness on the band of states near
maximal tension (\lean{D23d}), holds for every pair as well.

\textbf{The only thing a pair changes} is its floor,
\[
  \tfrac12|\langle[A,B]\rangle| \;=\; |a\,c|\cdot\tfrac12|\langle[T_d,P_d]\rangle| \;=\; \frac{|a\,c|}{d-1},
\]
expressed in the units of that pair.

\section{The table}

\begin{center}
\small
\renewcommand{\arraystretch}{1.25}
\begin{tabular}{rccccc}
\toprule
$d$ & $R=C_{\mathrm{Nava}}(d)$ & excess $\delta(d)$ & angle $\theta_{\mathrm{NRS}}(d)$
  & cube $\mathcal V(d,d,d)=\delta(d)^3$ & floor of a pair \\
\midrule
2 & 1 & 0 (saturates) & $0^\circ$ & 0 & $|ac|$ \\
3 & 1 & 0 (saturates) & $0^\circ$ & 0 & $|ac|/2$ \\
\textbf{4} & \textbf{1.008479} & \textbf{0.85\,\%} & $\mathbf{7.43^\circ}$
  & $\mathbf{6.09\times10^{-7}}$ & $|ac|/3$ \\
5 & 1.018350 & 1.84\,\% & $10.89^\circ$ & $6.18\times10^{-6}$ & $|ac|/4$ \\
6 & 1.027727 & 2.77\,\% & $13.34^\circ$ & $2.13\times10^{-5}$ & $|ac|/5$ \\
8 & 1.043563 & 4.36\,\% & $16.61^\circ$ & $8.27\times10^{-5}$ & $|ac|/7$ \\
10 & 1.055806 & 5.58\,\% & $18.71^\circ$ & $1.74\times10^{-4}$ & $|ac|/9$ \\
20 & 1.088392 & 8.84\,\% & $23.25^\circ$ & $6.91\times10^{-4}$ & $|ac|/19$ \\
50 & 1.114641 & 11.46\,\% & $26.21^\circ$ & $1.51\times10^{-3}$ & $|ac|/49$ \\
100 & 1.124785 & 12.48\,\% & $27.24^\circ$ & $1.94\times10^{-3}$ & $|ac|/99$ \\
\midrule
ceiling & $\sqrt{\pi^2/3-2}\approx1.135724$ & 13.57\,\% & $28.30^\circ$
  & $\delta_\infty^3\approx2.50\times10^{-3}$ & --- \\
\bottomrule
\end{tabular}
\end{center}

\textbf{How to read each column.}
\begin{itemize}
  \item $R=C_{\mathrm{Nava}}(d)$: the measured ratio $\sigma_A\sigma_B/|\langle\tilde A\psi^*,\tilde B\psi^*\rangle|$.
        It is the same for every row of the catalogue in Section~5.
  \item Excess $\delta(d)$: how far above the Robertson--Schrödinger floor the pair stays.
        Zero only at $d=2,3$; from $d=4$ on it is an algebraic quantum that no choice of
        units removes (\lean{D25}).
  \item Angle $\theta_{\mathrm{NRS}}(d)$: the opening between the fluctuation vectors of $A$
        and $B$. At $d=2,3$ they are parallel. From $d=4$ they open, at least $7.43^\circ$,
        and the opening grows with $d$ (\lean{D37b}).
  \item $\mathcal V(d,d,d)$: the volumetric quantum of the cube with three pairs, one per axis
        (\lean{D37e}). It is the box $\delta\times\delta\times\delta$. In the pictures of the
        repository the six fluctuation vectors of the three axes form a star whose envelope
        is an octahedron; each of its three axes opens by $\theta_{\mathrm{NRS}}(d)$.
  \item Floor of a pair: $\tfrac12|\langle[A,B]\rangle|$ in the units of that pair.
\end{itemize}

\textbf{The leftover does not disappear.} At $3\times3\times3$ nothing is left
($\mathcal V=0$, all three angles $0^\circ$). From $4\times4\times4$ something is always
left: $\mathcal V\ge\delta(4)^3$ (\lean{cuantoVolumetrico\_piso}) and every axis opens at least
$7.43^\circ$ (\lean{piso\_angular\_cubo}). Adding sites never shrinks it: $\mathcal V$ grows
strictly in each axis (\lean{cuantoVolumetrico\_strictMono\_x/y/z}) and the angle grows
strictly with $d$ (\lean{anguloNRS\_strictMonoOn}). It rises towards the ceiling
$\delta_\infty^3\approx2.50\times10^{-3}$ and $28.30^\circ$, and no cube reaches it
(\lean{cuantoVolumetrico\_techo}, \lean{anguloNRS\_lt\_limite}). The ceiling is a bound
approached by finite lattices, not a lattice of its own.

\section{The catalogue of pairs}

Each pair below is realized on $T_d{:}P_d$ with its own units $a,c$ and origins $b,e$; by the
theorem, all rows share the table of Section~4.

\begin{center}
\small
\renewcommand{\arraystretch}{1.2}
\begin{longtable}{r>{\raggedright}p{6.2cm}>{\raggedright}p{3.9cm}>{\raggedright\arraybackslash}p{3.9cm}}
\toprule
\# & Pair & $P_d$: the ordered values & $T_d$: the step \\
\midrule
\endhead
\multicolumn{4}{l}{\emph{Position and its equivalents}} \\
1 & position $x$ -- momentum $p_x$ & cell along $x$ & hop to the next cell \\
2 & position $y$ -- momentum $p_y$ & cell along $y$ & hop to the next cell \\
3 & position $z$ -- momentum $p_z$ & cell along $z$ & hop to the next cell \\
4 & centre of mass -- total momentum; relative coordinate -- relative momentum & cell of the coordinate & hop \\
5 & lattice site -- lattice momentum (tight-binding chain, open ends) & site & hop \\
6 & waveguide index -- transverse momentum & guide & evanescent coupling \\
7 & site of a spin chain, one-excitation sector -- exchange & excited site & exchange with the neighbour \\
8 & level $n$ of a $d$-level system -- drive between consecutive levels & level & drive $n\leftrightarrow n{+}1$ \\
\midrule
\multicolumn{4}{l}{\emph{Fields}} \\
9 & field $\varphi(x)$ -- conjugate momentum $\pi(x)$, site by site & field value & step of field value \\
10 & vector potential $A$ -- electric field $E$, mode by mode & value of $A$ & step of $A$ \\
11 & quadrature $X$ -- quadrature $Y$ of a mode & value of $X$ & step of $X$ \\
12 & photon number $n$ -- quadrature & photon number & $n\leftrightarrow n{+}1$ \\
13 & link electric field $E$ -- link variable $U$ (lattice gauge theory) & value of $E$ & $U$: $E\leftrightarrow E{+}1$ \\
\midrule
\multicolumn{4}{l}{\emph{Number and phase}} \\
14 & photon number $n$ -- phase $\varphi$ & photon number & phase shift $n\leftrightarrow n{+}1$ \\
15 & Cooper-pair number $n$ -- Josephson phase $\varphi$ & charge state & tunnelling $n\leftrightarrow n{+}1$ \\
16 & charge $Q$ -- flux $\Phi$ (LC circuit) & value of $Q$ & step of $Q$ \\
17 & number difference -- relative phase (two-mode condensate) & number difference & one particle between modes \\
\midrule
\multicolumn{4}{l}{\emph{Angle, rotation and spin}} \\
18 & angle $\varphi$ -- angular momentum $L_z$ & value $m$ of $L_z$ & $m\leftrightarrow m{+}1$ \\
19 & angle -- action (action--angle variables) & value of the action & step of the action \\
20 & clock $Z$ -- shift $X$ (Weyl pair, $d$ levels; \lean{D12}) & clock value & shift by one \\
21 & $\sigma_z$ -- $\sigma_x$ of a qubit & $\pm1$ & flip; exactly $(P_2,T_2)$ \\
22 & $J_z$ -- $J_x$ (and cyclic), spin or angular momentum & value $m$ of $J_z$ & $m\leftrightarrow m{+}1$ \\
23 & isospin and colour generators & weight along one generator & step between weights \\
\midrule
\multicolumn{4}{l}{\emph{Magnetic field}} \\
24 & kinetic momenta $\pi_x$ -- $\pi_y$ (Landau levels) & Landau level & $n\leftrightarrow n{+}1$ \\
25 & guiding centre $X$ -- $Y$ & guiding-centre level & $n\leftrightarrow n{+}1$ \\
\midrule
\multicolumn{4}{l}{\emph{Time}} \\
26 & time $t$ -- energy $E$ & registered tick & step between ticks \\
27 & time -- frequency of a signal (Fourier, Gabor) & time bin & step between bins \\
\bottomrule
\end{longtable}
\end{center}

Rows 26 and 27: in the theorem of 1929--1930 time is not an operator, and the reading of time
as registered ticks belongs to the physical layer, outside this repository. Row 27 concerns
classical signals with the same algebra as row~1.

Rows 1--3 are the three axes of the cube of \lean{D37}. Any three rows whose pairs commute with
each other can occupy the three axes, and items 7--9 of the theorem apply to them.

\section{Reproduce}

\begin{verbatim}
lake exe cache get
lake build                                       # the whole package
lake env lean Verification/Layer1_Mathematics.lean  # boundary + axioms
\end{verbatim}
Every theorem cited here depends only on \lean{propext}, \lean{Classical.choice} and
\lean{Quot.sound}; the package has no \lean{sorry}. The numbers of the table are
$C_{\mathrm{Nava}}(d)$ from \lean{D9}/\lean{D21}, recomputed from the matrices $T_d$, $P_d$ and
the vector $\psi^*$; the exact value at $d=4$ is
$C_{\mathrm{Nava}}(4)=\sqrt{(99-42\sqrt5)/5}$ (\lean{D37b}, \lean{D37e}).

\end{document}
